\documentclass[11pt]{article}
\usepackage[margin=1in]{geometry}
\usepackage{booktabs}
\usepackage{graphicx}
\usepackage{hyperref}
\usepackage{enumitem}
\setlength{\parindent}{0pt}
\setlength{\parskip}{0.6em}

\title{Prompt-Profile Binary Capacity Controls On Balanced \texttt{majority\_s\_0.5}}
\author{}
\date{2026-04-04 06:31 UTC}

\begin{document}
\maketitle

\section*{Executive Summary}
\begin{itemize}[leftmargin=1.5em]
\item This note compares three probe families on the same saved April binary data: default \texttt{h128 d1}, width-only control \texttt{h256 d1}, and width+depth control \texttt{h256 d2}.
\item Exact runtime:
  \begin{itemize}[leftmargin=1.5em]
  \item \texttt{2106}: \texttt{LiveCodeBench}-only pilot caused by the first dataset-list bug
  \item \texttt{2107}: corrected five-dataset \texttt{h256 d2} rerun, \texttt{2026-04-04 05:57:19 UTC} to \texttt{2026-04-04 06:04:55 UTC}
  \item \texttt{2108}: width-only \texttt{h256 d1} control, \texttt{2026-04-04 06:13:02 UTC} to \texttt{2026-04-04 06:22:07 UTC}
  \end{itemize}
\item Fixed object: same saved prompt-prefill activations, same binary target \texttt{majority\_s\_0.5}, same train-balanced / test-natural split contract, same seeds \texttt{0,1,2}, same \texttt{epochs=15}, \texttt{batch\_size=256}, \texttt{lr=1e-4}, \texttt{weight\_decay=0.05}, \texttt{dropout=0.1}.
\item Main read on held-out \texttt{PR-AUC}: for the per-layer \texttt{ensemble}, width is the useful change and added depth hurts. For \texttt{last\_layer}, added depth helps modestly on top of width.
\item Cross-dataset mean test \texttt{PR-AUC}: prompt length baseline \texttt{0.350}; ensemble \texttt{h128 d1} \texttt{0.474}; ensemble \texttt{h256 d1} \texttt{0.518}; ensemble \texttt{h256 d2} \texttt{0.492}; last-layer \texttt{h128 d1} \texttt{0.410}; last-layer \texttt{h256 d1} \texttt{0.420}; last-layer \texttt{h256 d2} \texttt{0.430}.
\item Best single global surface from this control sweep: \texttt{ensemble h256 d1}.
\item Secondary sanity check: under \texttt{best\_rank}, the ensemble ordering stays the same (\texttt{h256 d1 0.539 > h256 d2 0.522 > h128 d1 0.474}), while the \texttt{last\_layer} ordering does not. That means the robust tuning conclusion is the ensemble one.
\item Threshold caveat: on the natural test split, rare-positive datasets still show recall-heavy operating points after train balancing. Keep \texttt{PR-AUC} primary and treat threshold metrics as behavior diagnostics only.
\item Limitation: this still does not answer which layers to keep. It only separates width from depth while keeping the same two view definitions: \texttt{ensemble} over all \texttt{28} layers, or final \texttt{last\_layer}.
\end{itemize}

\section*{Exact Question}
On the saved balanced binary object, did the earlier \texttt{h256 d2} shifts come from width, depth, or only the combination?

\section*{Run Contract}
\begin{itemize}[leftmargin=1.5em]
\item Source data root: \path{/data/scratch/murphy/outputs/cot-loop-detection/full_train_locked_pair_20260404}
\item Binary target: \texttt{majority\_s\_0.5}
\item Train balance: downsample to \texttt{50/50}
\item Test balance: natural
\item Feature surface: stacked prompt-prefill activations with sample shape \texttt{[28, 2048]}
\item Views:
  \begin{itemize}[leftmargin=1.5em]
  \item \texttt{ensemble}: one MLP per layer across all \texttt{28} saved layers, aggregated by \texttt{vote\_fraction}
  \item \texttt{last\_layer}: final-layer-only MLP
  \end{itemize}
\item Probe families:
  \begin{itemize}[leftmargin=1.5em]
  \item \texttt{h128 d1}: \texttt{hidden\_dim=128}, \texttt{depth=1}
  \item \texttt{h256 d1}: \texttt{hidden\_dim=256}, \texttt{depth=1}
  \item \texttt{h256 d2}: \texttt{hidden\_dim=256}, \texttt{depth=2}
  \end{itemize}
\item Optimization held fixed: \texttt{epochs=15}, \texttt{batch\_size=256}, \texttt{lr=1e-4}, \texttt{weight\_decay=0.05}, \texttt{dropout=0.1}
\end{itemize}

\section*{Ranking Results: Test \texttt{PR-AUC}}
\begin{center}
\small
\setlength{\tabcolsep}{4pt}
\resizebox{\textwidth}{!}{
\begin{tabular}{@{}lrrrrrrrr@{}}
\toprule
Dataset & Test prevalence & Prompt length & Ensemble \texttt{h128 d1} & Ensemble \texttt{h256 d1} & Ensemble \texttt{h256 d2} & Last-layer \texttt{h128 d1} & Last-layer \texttt{h256 d1} & Last-layer \texttt{h256 d2} \\
\midrule
GPQA & 0.0500 & 0.0657 & 0.4198 & 0.5833 & 0.5025 & 0.2298 & 0.3071 & 0.3176 \\
AIME & 0.5833 & 0.8976 & 0.9218 & 0.9120 & 0.9090 & 0.9043 & 0.8436 & 0.8479 \\
MATH-500 & 0.0400 & 0.1002 & 0.1354 & 0.1663 & 0.1596 & 0.1506 & 0.1315 & 0.1322 \\
MMLU-Pro & 0.0125 & 0.1104 & 0.1837 & 0.2116 & 0.2023 & 0.0559 & 0.0701 & 0.0959 \\
LiveCodeBench & 0.3375 & 0.5760 & 0.7111 & 0.7143 & 0.6865 & 0.7116 & 0.7452 & 0.7575 \\
\bottomrule
\end{tabular}
}
\end{center}

\section*{Cross-Dataset Mean \texttt{PR-AUC}}
\begin{center}
\small
\begin{tabular}{@{}lr@{}}
\toprule
Surface & Mean test \texttt{PR-AUC} \\
\midrule
Prompt length baseline & 0.3500 \\
Ensemble \texttt{h128 d1} & 0.4744 \\
Ensemble \texttt{h256 d1} & 0.5175 \\
Ensemble \texttt{h256 d2} & 0.4920 \\
Last-layer \texttt{h128 d1} & 0.4104 \\
Last-layer \texttt{h256 d1} & 0.4195 \\
Last-layer \texttt{h256 d2} & 0.4302 \\
\bottomrule
\end{tabular}
\end{center}

\section*{Secondary Check: \texttt{best\_rank} Mean Test \texttt{PR-AUC}}
\begin{center}
\small
\begin{tabular}{@{}lr@{}}
\toprule
Surface & Mean test \texttt{PR-AUC} \\
\midrule
Prompt length baseline & 0.3500 \\
Ensemble \texttt{h128 d1} & 0.4744 \\
Ensemble \texttt{h256 d1} & 0.5385 \\
Ensemble \texttt{h256 d2} & 0.5215 \\
Last-layer \texttt{h128 d1} & 0.4104 \\
Last-layer \texttt{h256 d1} & 0.4552 \\
Last-layer \texttt{h256 d2} & 0.4359 \\
\bottomrule
\end{tabular}
\end{center}

This secondary checkpoint rule keeps the global ensemble recommendation intact: \texttt{h256 d1} still stays ahead of \texttt{h256 d2}. The \texttt{last\_layer} ranking does move, which is another reason not to treat that control view as the main tuning conclusion.

\section*{Threshold Behavior On Natural Test For The Recommended Surface}
Held fixed to \texttt{ensemble h256 d1} at the frozen \texttt{best\_loss} checkpoint.

\begin{center}
\small
\setlength{\tabcolsep}{4pt}
\begin{tabular}{@{}lrrrrr@{}}
\toprule
Dataset & Test prevalence & Positive precision & Positive recall & Positive \texttt{F1} & Macro \texttt{F1} \\
\midrule
GPQA & 0.0500 & 0.2143 & 1.0000 & 0.3452 & 0.6069 \\
AIME & 0.5833 & 0.9048 & 0.4762 & 0.5639 & 0.6143 \\
MATH-500 & 0.0400 & 0.0539 & 0.6667 & 0.0996 & 0.4391 \\
MMLU-Pro & 0.0125 & 0.1249 & 1.0000 & 0.2176 & 0.5746 \\
LiveCodeBench & 0.3375 & 0.6143 & 0.6543 & 0.5138 & 0.5809 \\
\bottomrule
\end{tabular}
\end{center}

Because training is balanced but evaluation is natural, these threshold metrics stay recall-heavy on rare-positive datasets. That is expected behavior on this object, not evidence that the operating point is suddenly well-calibrated. Use them as a threshold-shape check, not as the headline comparison.

\section*{What Width And Depth Actually Did}
\begin{itemize}[leftmargin=1.5em]
\item Ensemble:
  \begin{itemize}[leftmargin=1.5em]
  \item width helps relative to the default on \texttt{GPQA}, \texttt{MATH-500}, \texttt{MMLU-Pro}, and \texttt{LiveCodeBench}, with only a small drop on \texttt{AIME}
  \item added depth hurts relative to width-only on all five ensemble slices
  \item so the best global ensemble surface is \texttt{h256 d1}, not \texttt{h256 d2}
  \end{itemize}
\item Last-layer:
  \begin{itemize}[leftmargin=1.5em]
  \item width helps on \texttt{GPQA}, \texttt{MMLU-Pro}, and \texttt{LiveCodeBench}, but hurts on \texttt{AIME} and \texttt{MATH-500}
  \item added depth then improves on top of width on all five last-layer slices, though sometimes only slightly
  \item so the best global last-layer surface is \texttt{h256 d2}
  \end{itemize}
\end{itemize}

\section*{Interpretation}
\begin{itemize}[leftmargin=1.5em]
\item Proven now:
  \begin{itemize}[leftmargin=1.5em]
  \item \texttt{ensemble h256 d1} is the strongest single global surface in this capacity sweep under the frozen \texttt{best\_loss} rule
  \item \texttt{ensemble h256 d1} also stays ahead of \texttt{h256 d2} under the secondary \texttt{best\_rank} rule, so the main ensemble choice is not a checkpoint-selection artifact
  \item \texttt{ensemble h256 d1} still beats the prompt-length baseline on all five datasets by \texttt{PR-AUC}
  \item \texttt{last\_layer h256 d2} is the best last-layer variant only under the frozen \texttt{best\_loss} rule, and even there it is still weaker than \texttt{ensemble h256 d1} as a global choice
  \end{itemize}
\item Also proven:
  \begin{itemize}[leftmargin=1.5em]
  \item the earlier \texttt{h256 d2} ensemble result was mixing one good change with one bad one
  \item the good change was width
  \item the bad change, for ensemble, was the extra depth
  \end{itemize}
\item Important nuance:
  \begin{itemize}[leftmargin=1.5em]
  \item the \texttt{last\_layer} ordering is not stable across checkpoint rules (\texttt{h256 d2} wins at \texttt{best\_loss}, \texttt{h256 d1} wins at \texttt{best\_rank})
  \item so the robust tuning conclusion is about the ensemble, not about promoting a separate last-layer champion
  \end{itemize}
\item Important caveat:
  \begin{itemize}[leftmargin=1.5em]
  \item threshold metrics on the natural test split remain recall-heavy on rare-positive datasets because train balance and test prevalence differ
  \item keep \texttt{PR-AUC} primary and treat threshold metrics as diagnostic only
  \item this still says nothing about whether we should keep all \texttt{28} layers in the ensemble
  \item the next honest tuning step is a small layer-subset / view sweep on the same balanced binary data, not another vague capacity increase
  \end{itemize}
\end{itemize}

\section*{Recommendation}
\begin{itemize}[leftmargin=1.5em]
\item If one single global binary surface is needed today, use \texttt{ensemble h256 d1}.
\item Keep \texttt{PR-AUC} as the main comparison object for this balanced-train / natural-test binary surface.
\item Keep \texttt{best\_loss} as the frozen main checkpoint rule and \texttt{best\_rank} as a sanity check, not as a reason to reopen the global ensemble choice.
\item If the note still needs a cheap \texttt{last\_layer} control under the same frozen rule, use \texttt{h256 d2}, but do not present that as a stable tuning winner.
\end{itemize}

\section*{Artifact Bundle}
\begin{itemize}[leftmargin=1.5em]
\item PDF: \path{outputs/prompt_profile_binary_capacity_controls_20260404/prompt_profile_binary_capacity_controls_20260404.pdf}
\item TeX: \path{outputs/prompt_profile_binary_capacity_controls_20260404/prompt_profile_binary_capacity_controls_20260404.tex}
\item Summary JSON: \path{outputs/prompt_profile_binary_capacity_controls_20260404/capacity_comparison_summary.json}
\item Summary CSV: \path{outputs/prompt_profile_binary_capacity_controls_20260404/capacity_comparison_metrics.csv}
\item Intermediate \texttt{h256 d2} note:
  \begin{itemize}[leftmargin=1.5em]
  \item \texttt{docs/prompt-profile-binary-retrain-h256d2-2026-04-04.md}
  \item use this only as the raw \texttt{2106} / \texttt{2107} rerun record, not as the current recommendation surface
  \end{itemize}
\end{itemize}

\end{document}
